\hypertarget{cont7:paper-ix-weighted-first-order-root-motion}{%
\subsection{Chapter IX: weighted first-order root
motion}\label{cont7:paper-ix-weighted-first-order-root-motion}}

\hypertarget{cont7:the-theorem}{%
\subsubsection{The theorem}\label{cont7:the-theorem}}

Let \(p\) be prime, \(R=\mathbb Z_p\), \(K=\mathbb Q_p\), and
\(d,s,M\in\mathbb N\). Let \(F,H\in R[X_1,\ldots,X_d]^d\), let
\(a\in R^d\) satisfy \(F(a)=0\), and set \[
J=DF(a),\qquad G=F+p^M H.
\] Assume \[
J\in\operatorname{GL}_d(K),\qquad p^sJ^{-1}\in M_d(R),\qquad M>2s.
\] Then there is exactly one root \(b\) of \(G\) in \(a+p^{s+1}R^d\). It
satisfies \[
b-a\in p^{M-s}R^d,
\] and there is \(\varepsilon\in p^{2M-3s}R^d\) such that, in \(K^d\),
\[
\boxed{b-a=-p^M J^{-1}H(a)+\varepsilon.}
\] The inverse in this formula belongs to the original Jacobian
\(DF(a)\). The theorem constructs the perturbed root from these
hypotheses and proves uniqueness throughout the full ball of depth
\(s+1\). It does not require an additional Jacobian hypothesis on \(G\).
Coefficients and centers are arbitrary integral p-adic elements. Every
prime and finite square dimension are included, including \(p=2\) and
dimension zero.

The error depth exceeds the displacement depth by \[
(2M-3s)-(M-s)=M-2s\ge1.
\] Thus the first-order term determines strictly more digits than the
displacement bound alone. These are guaranteed divisibility bounds;
equality of valuations is not asserted for every system.

\hypertarget{cont7:root-construction-with-the-original-inverse}{%
\subsubsection{Root construction with the original
inverse}\label{cont7:root-construction-with-the-original-inverse}}

Write \(A=p^sJ^{-1}\in M_d(R)\), so \(AJ=JA=p^sI\). Let \(t=p^{s+1}\)
and \(k=M-(2s+1)\ge0\). The actual polynomial Taylor and difference
identities give polynomial systems \(Q_F,Q_H\), both vanishing at zero,
such that \[
(AF)(a+tu)=p^{2s+1}u+p^{2s+2}Q_F(u),
\] \[
(AH)(a+tu)=(AH)(a)+tQ_H(u).
\] Consequently \[
(AG)(a+tu)=p^{2s+1}\bigl(C+u+pQ(u)\bigr),
\] where \[
C=p^kAH(a),\qquad
Q=Q_F+p^{M-s-1}Q_H,\qquad Q(0)=0.
\] All exponents are nonnegative under \(M>2s\). Actual integral
polynomial evaluation is 1-Lipschitz, so \(u\mapsto-C-pQ(u)\) is a
strict contraction on the complete space \(R^d\). Its unique fixed point
inherits the depth \(k\) of \(C\), because \(Q(0)=0\). Thus \(b=a+tu\)
has displacement depth \(s+1+k=M-s\).

The certificate \(JA=p^sI\) makes multiplication by \(A\) injective, so
the normalized zero is an actual root of \(G\). Every root in \(a+tR^d\)
has the form \(a+tw\) and satisfies the same fixed-point equation. This
proves uniqueness on the entire stated ball using the original matrix
certificate.

\hypertarget{cont7:the-first-order-error}{%
\subsubsection{The first-order error}\label{cont7:the-first-order-error}}

Put \(\delta=b-a\) and \(L=M-s\). The displacement form of Taylor's
theorem, proved over any commutative ring, gives \[
E_F=(AF)(b)-(AF)(a)-p^s\delta\in p^{2L}R^d.
\] Polynomial congruence gives \[
E_H=(AH)(b)-(AH)(a)\in p^LR^d.
\] Since \(F(a)=0\) and \(G(b)=0\), \[
p^s\delta+p^MAH(a)=-E_F-p^ME_H.
\] We have \(M+L\ge2L\), so the right side lies in \(p^{2L}R^d\).
Cancelling the nonzero scalar \(p^s\) proves the integral estimate \[
\delta+p^{M-s}AH(a)\in p^{2M-3s}R^d.
\] Define this integral vector to be \(\varepsilon\). Under
\(R\hookrightarrow K\), the exact identity \(p^{M-s}A=p^MJ^{-1}\) gives
the displayed first-order formula. The Lean bridge checks matrix
evaluation and all coefficient inclusions explicitly.

\hypertarget{cont7:sharpness-and-the-strict-threshold}{%
\subsubsection{Sharpness and the strict
threshold}\label{cont7:sharpness-and-the-strict-threshold}}

For \(p=2\), take \(F(X)=X+X^2\), \(H(X)=1\), \(a=0\), \(s=0\), \(M=1\).
There is a unique root \(b\in2\mathbb Z_2\) of \(X+X^2+2\). Its equation
gives \(b+2=-b^2\). It also forces \(b\equiv2\pmod4\). Therefore \[
v_2(b+2)=2=2M-3s.
\] The accompanying Lean controls show that the error is divisible by
four and is not divisible by eight. This refutes a uniform replacement
of \(2M-3s\) by \(2M-3s+1\), even for a one-variable unit Jacobian.

The prior strict-threshold example also applies to perturbations: take
\(F(X)=X^2-1\), \(H(X)=1\), \(a=1\), \(p=2\), \(s=1\), \(M=2\). Then
\(F(a)=0\), \(DF(a)=2\), and \(F+2^2H=X^2+3\) has no root modulo eight,
hence none in \(\mathbb Z_2\). The new explanatory identification uses
the previously checked no-root polynomial. The strict inequality
\(M>2s\) cannot uniformly be replaced by \(M\ge2s\).
